
        \documentclass[spanish, 11pt]{exam}

        %These tell TeX which packages to use.
        \usepackage{array,epsfig}
        \usepackage{amsmath, textcomp}
        \usepackage{amsfonts}
        \usepackage{amssymb}
        \usepackage{amsxtra}
        \usepackage{amsthm}
        \usepackage{mathrsfs}
        \usepackage{color}
        \usepackage{multicol, xparse}
        \usepackage{verbatim}


        \usepackage[utf8]{inputenc}
        \usepackage[spanish]{babel}
        \usepackage{eurosym}

        \usepackage{graphicx}
        \graphicspath{{../img/}}



        \printanswers
        \nopointsinmargin
        \pointformat{}

        %Pagination stuff.
        %\setlength{\topmargin}{-.3 in}
        %\setlength{\oddsidemargin}{0in}
        %\setlength{\evensidemargin}{0in}
        %\setlength{\textheight}{9.in}
        %\setlength{\textwidth}{6.5in}
        %\pagestyle{empty}

        \let\multicolmulticols\multicols
        \let\endmulticolmulticols\endmulticols
        \RenewDocumentEnvironment{multicols}{mO{}}
         {%
          \ifnum#1=1
            #2%
          \else % More than 1 column
            \multicolmulticols{#1}[#2]
          \fi
         }
         {%
          \ifnum#1=1
          \else % More than 1 column
            \endmulticolmulticols
          \fi
         }
        \renewcommand{\solutiontitle}{\noindent\textbf{Sol:}\enspace}

        \newcommand{\samedir}{\mathbin{\!/\mkern-5mu/\!}}

        \newcommand{\class}{1º Bachillerato}
        \newcommand{\examdate}{\today}

        \newcommand{\tipo}{A}


        \newcommand{\timelimit}{50 minutos}



        \pagestyle{head}
        \firstpageheader{\includegraphics[width=0.2\columnwidth]{header_left}}{\textbf{Departamento de Matemáticas\linebreak \class}\linebreak \examnum}{\includegraphics[width=0.1\columnwidth]{header_right}}
        \runningheader{\class}{\examnum}{Página \thepage\ of \numpages}
        \runningheadrule

        \newcommand{\examnum}{9 - Ecuaciones exponenciales y logarítmicas}
        \begin{document}
        \begin{questions}
        \question p027e04 - Resuelve las siguientes ecuaciones exponenciales:
        \begin{multicols}{2} 
        \begin{parts} \part[1]  $ {10^{3 - x}} = 1 $  \begin{solution}  $ \left [ 3\right ] $  \end{solution} \part[1]  $ {5^{x + 3}} = 125 $  \begin{solution}  $ \left [ 0\right ] $  \end{solution} \part[1]  $ 5^{1-x^2}=\frac{1}{125} $  \begin{solution}  $ \left [ -2, \quad 2\right ] $  \end{solution} \part[1]  $ {5^{{x^2} - 5x + 6}} = 1 $  \begin{solution}  $ \left [ 2, \quad 3\right ] $  \end{solution} \part[1]  $ {2^{1 - x}} = \frac{1}{8} $  \begin{solution}  $ \left [ 4\right ] $  \end{solution} \part[1]  $ 2^{x + 3} = 4^{- x} $  \begin{solution}  $ \left [ -1\right ] $  \end{solution} \part[1]  $ {9^{x - 1}} = {3^{x + 1}} $  \begin{solution}  $ \left [ 3\right ] $  \end{solution} \part[1]  $ {4^{4x + 3}} = {2^{ - x}} $  \begin{solution}  $ \left [ - \frac{2}{3}\right ] $  \end{solution} \part[1]  $ {8^{x - 1}} = {4^{3x + 1}} $  \begin{solution}  $ \left [ - \frac{5}{3}\right ] $  \end{solution} \part[1]  $ {5^{ - x}} = 0.04 $  \begin{solution}  $ \left [ 2\right ] $  \end{solution} \part[1]  $ {10^x} \cdot {10^{ - 2x + 7}} = 100 $  \begin{solution}  $ \left [ 5\right ] $  \end{solution} \part[1]  $ {( {{3^x}} )^2} \cdot {3^x} = {9^3} $  \begin{solution}  $ \left [ 2\right ] $  \end{solution} \part[1]  $ \sqrt{2 \cdot \sqrt{2 \cdot \sqrt{2} } } = {2^x} $  \begin{solution}  $ \left [ \frac{7}{8}\right ] $  \end{solution} \part[1]  $ {2^{{x^2} - 5x}} = {64^{ - 1}} $  \begin{solution}  $ \left [ 2, \quad 3\right ] $  \end{solution} \part[1]  $ \sqrt{\sqrt {3}  + \sqrt {3}  + \sqrt {3} }  = {3^{x + 2}} $  \begin{solution}  $ \left [ - \frac{5}{4}\right ] $  \end{solution} \part[1]  $ \sqrt[x]{216}  = 6 $  \begin{solution}  $ \left [ 3\right ] $  \end{solution} \part[1]  $ {4^x} - {2^x} = 2 $  \begin{solution}  $ \left [ 1\right ] $  \end{solution} \part[1]  $ {5^x} - 30 \cdot {5^x} + 145 = 0 $  \begin{solution}  $ \left [ 1\right ] $  \end{solution} \part[1]  $ {2^{x - 1}} + {2^x} + {2^{x + 1}} = 7 $  \begin{solution}  $ \left [ 1\right ] $  \end{solution}
        \end{parts}
        \end{multicols}
        \question p028e05 - Resuelve las siguientes ecuaciones exponenciales:
        \begin{multicols}{2} 
        \begin{parts} \part[1]  $ {3^{x + 1}} + {3^x} + {3^{x - 1}} = 117 $  \begin{solution}  $ \left [ 3\right ] $  \end{solution} \part[1]  $ {3^x} + {3^{x - 1}} + {3^{x - 2}} + {3^{x - 3}} + {3^{x - 4}} = 363 $  \begin{solution}  $ \left [ 5\right ] $  \end{solution} \part[1]  $ {2^{3x}} - \frac{3}{{{2^{3x + 2}}}} + 1 = 0 $  \begin{solution}  $ \left [ - \frac{1}{3}\right ] $  \end{solution} \part[1]  $ {3^{x - 1}} + {3^{2 - x}} = 4 $  \begin{solution}  $ \left [ 1, \quad 2\right ] $  \end{solution} \part[1]  $ {2^{x + 1}} + {4^x} = 80 $  \begin{solution}  $ \left [ 3\right ] $  \end{solution} \part[1]  $ 2^{2x} - 3 \cdot {2^{x + 1}} + 8 = 0 $  \begin{solution}  $ \left [ 1, \quad 2\right ] $  \end{solution} \part[1]  $ {3^{2x - 3}} + 1 = 4 \cdot {3^{x - 2}} $  \begin{solution}  $ \left [ 1, \quad 2\right ] $  \end{solution} \part[1]  $ {2^{2x}} - 10 \cdot {2^x} + 16 = 0 $  \begin{solution}  $ \left [ 1, \quad 3\right ] $  \end{solution} \part[1]  $ {16^x} - {4^x} = 240 $  \begin{solution}  $ \left [ 2\right ] $  \end{solution} \part[1]  $ 9^x - 6 \cdot {3^{x + 1}} + 81 = 0 $  \begin{solution}  $ \left [ 2\right ] $  \end{solution} \part[1]  $ {3^{x + 2}} + {9^{x + 1}} = 810 $  \begin{solution}  $ \left [ 2\right ] $  \end{solution} \part[1]  $ {5^{x - 1}} = 2 + \frac{3}{{{5^{x - 2}}}} $  \begin{solution}  $ \left [ 2\right ] $  \end{solution} \part[1]  $ {3^{x + 1}} + {3^{x - 2}} = \frac{{15}}{{{3^{x - 1}}}} + \frac{{247}}{{{3^{x - 2}}}} $  \begin{solution}  $ \left [ 3\right ] $  \end{solution} \part[1]  $ {4^{2x}} + 16 \cdot {4^{ - 2x}} - 10 = 0 $  \begin{solution}  $ \left [ \frac{1}{4}, \quad \frac{3}{4}\right ] $  \end{solution}
        \end{parts}
        \end{multicols}
        \question p028e06 - Resuelve los siguientes sistemas:
        \begin{multicols}{2} 
        \begin{parts} \part[1]  $ \left\{\begin{matrix}{3^x} = {3^y}\\{4^x} \cdot {4^y} = 256\end{matrix}\right. $  \begin{solution}  $ \left [ \left \{ x : 2, \quad y : 2\right \}\right ] $  \end{solution} \part[1]  $ \left\{\begin{matrix}{2^{x + 2y}} = 32\\{2^{3x - 5y}} = 16\end{matrix}\right. $  \begin{solution}  $ \left [ \left \{ x : 3, \quad y : 1\right \}\right ] $  \end{solution} \part[1]  $ \left\{\begin{matrix}{5^x} = {5^y} \cdot 625 \\{2^x} \cdot {2^y} = 256\end{matrix}\right. $  \begin{solution}  $ \left [ \left \{ x : 6, \quad y : 2\right \}\right ] $  \end{solution} \part[1]  $ \left\{\begin{matrix}{2^x} + {2^y} = 24\\{2^{x + y}} = 128\end{matrix}\right. $  \begin{solution}  $ \left [ \left \{ x : 3, \quad y : 4\right \}, \quad \left \{ x : 4, \quad y : 3\right \}\right ] $  \end{solution}
        \end{parts}
        \end{multicols}
        \question p028e07 - Calcula:
        \begin{multicols}{4} 
        \begin{parts} \part[1]  $ \log 100 $  \begin{solution}  $ 2 $  \end{solution} \part[1]  $ \log_{5}(625) $  \begin{solution}  $ 4 $  \end{solution} \part[1]  $ \log_{2}(32) $  \begin{solution}  $ 5 $  \end{solution} \part[1]  $ \log_{3}(81) $  \begin{solution}  $ 4 $  \end{solution} \part[1]  $ \log_{2}(1024) $  \begin{solution}  $ 10 $  \end{solution} \part[1]  $ \log 1000 $  \begin{solution}  $ 3 $  \end{solution} \part[1]  $ \log 10000 $  \begin{solution}  $ 4 $  \end{solution} \part[1]  $ \log 1000000 $  \begin{solution}  $ 6 $  \end{solution} \part[1]  $ \log (10^6) $  \begin{solution}  $ 6 $  \end{solution} \part[1]  $ \log 0.1 $  \begin{solution}  $ -1 $  \end{solution} \part[1]  $ \log 0.01 $  \begin{solution}  $ -2 $  \end{solution} \part[1]  $ \log 0.001 $  \begin{solution}  $ -3 $  \end{solution} \part[1]  $ \log 0.000001 $  \begin{solution}  $ -6 $  \end{solution} \part[1]  $ \log_{5}(625) $  \begin{solution}  $ 4 $  \end{solution} \part[1]  $ \log_{2}(4) $  \begin{solution}  $ 2 $  \end{solution} \part[1]  $ \log_{2}(64) $  \begin{solution}  $ 6 $  \end{solution}
        \end{parts}
        \end{multicols}
        \question p028e07b - Calcula (continuación):
        \begin{multicols}{4} 
        \begin{parts} \part[1]  $ \log_{2}(\frac{1}{2}) $  \begin{solution}  $ -1 $  \end{solution} \part[1]  $ \log_{2}(\frac{1}{4}) $  \begin{solution}  $ -2 $  \end{solution} \part[1]  $ \log_{2}(\sqrt{2}) $  \begin{solution}  $ \frac{1}{2} $  \end{solution} \part[1]  $ \log_{2}(\sqrt{8}) $  \begin{solution}  $ \frac{3}{2} $  \end{solution} \part[1]  $ \log_{3}(3) $  \begin{solution}  $ 1 $  \end{solution} \part[1]  $ \log_{3}(27) $  \begin{solution}  $ 3 $  \end{solution} \part[1]  $ \log_{3}(27) $  \begin{solution}  $ 3 $  \end{solution} \part[1]  $ \log_{3}(\frac{1}{3}) $  \begin{solution}  $ -1 $  \end{solution} \part[1]  $ \log_{3}(\frac{1}{9}) $  \begin{solution}  $ -2 $  \end{solution} \part[1]  $ \log_{3}{\sqrt[3]{3}} $  \begin{solution}  $ \frac{1}{3} $  \end{solution} \part[1]  $ \log_{\frac{1}{3}}(81) $  \begin{solution}  $ -4 $  \end{solution} \part[1]  $ \log_{0.8}(1) $  \begin{solution}  $ 0 $  \end{solution} \part[1]  $ \log_{0.01}({10}^{-3}) $  \begin{solution}  $ \frac{3}{2} $  \end{solution} \part[1]  $ \log_{\frac{1}{49}}(7) $  \begin{solution}  $ - \frac{1}{2} $  \end{solution} \part[1]  $ \log_{\frac{1}{5}}({\frac{1}{25}})^{\frac{1}{5}} $  \begin{solution}  $ \sqrt[5]{2} $  \end{solution}
        \end{parts}
        \end{multicols}
        \question p028e08 - Averigua el valor de x en los siguientes casos:
        \begin{multicols}{2} 
        \begin{parts} \part[1]  $ \log_{7}(x) = 2 $  \begin{solution}  $ \left [ 49\right ] $  \end{solution} \part[1]  $ \log_{8}(x) = \frac{1}{3} $  \begin{solution}  $ \left [ 2\right ] $  \end{solution} \part[1]  $ \log_{2}(\frac{1}{16}) = x $  \begin{solution}  $ \left [ -4\right ] $  \end{solution} \part[1]  $ \log_{x}(125) = 3 $  \begin{solution}  $ \left [ 5\right ] $  \end{solution} \part[1]  $ \log_{2}(64) = x $  \begin{solution}  $ \left [ 6\right ] $  \end{solution} \part[1]  $ \log_{x}(9) = 2 $  \begin{solution}  $ \left [ 3\right ] $  \end{solution} \part[1]  $ \log_{2}(x) = -3 $  \begin{solution}  $ \left [ \frac{1}{8}\right ] $  \end{solution} \part[1]  $ \log_{x}(5) + 2 $  \begin{solution}  $ \left [ \frac{\sqrt{5}}{5}\right ] $  \end{solution} \part[1]  $ \log_{0.008}(625) = 2x $  \begin{solution}  $ \left [ - \frac{2}{3}\right ] $  \end{solution}
        \end{parts}
        \end{multicols}
        \question p028e09 - Sabiendo que log 2 = 0,301030, calcula
        \begin{multicols}{3} 
        \begin{parts} \part[1]  $ \log(16) $  \begin{solution}  $ 1.20411998 $  \end{solution} \part[1]  $ \log(64) $  \begin{solution}  $ 1.80617997 $  \end{solution} \part[1]  $ \log(1024) $  \begin{solution}  $ 3.01029996 $  \end{solution} \part[1]  $ \log(\frac{1}{2}) $  \begin{solution}  $ -0.301029996 $  \end{solution} \part[1]  $ \log(\frac{1}{16}) $  \begin{solution}  $ -1.20411998 $  \end{solution} \part[1]  $ \log(5) $  \begin{solution}  $ 0.698970004 $  \end{solution} \part[1]  $ \log(25) $  \begin{solution}  $ 1.39794001 $  \end{solution} \part[1]  $ \log(0.0016) $  \begin{solution}  $ -2.79588002 $  \end{solution} \part[1]  $ \log(\sqrt[ 4 ]{\frac{1}{{0.04}}} ) $  \begin{solution}  $ 0.349485002 $  \end{solution} \part[1]  $ \log{\sqrt[ 3 ]{0.002}} $  \begin{solution}  $ -0.899656668 $  \end{solution} \part[1]  $ \log(0.025) $  \begin{solution}  $ -1.60205999 $  \end{solution} \part[1]  $ \log(\sqrt[ 4 ]{\frac{1}{{1024}}} ) $  \begin{solution}  $ -0.752574989 $  \end{solution}
        \end{parts}
        \end{multicols}
        \question p028e10 - Sabiendo que log 3 = 0,477121, calcula
        \begin{multicols}{3} 
        \begin{parts} \part[1]  $ \log(243) $  \begin{solution}  $ 2.38560627 $  \end{solution} \part[1]  $ \log(0.0027) $  \begin{solution}  $ -2.56863624 $  \end{solution} \part[1]  $ \log(0.003) $  \begin{solution}  $ -2.52287875 $  \end{solution} \part[1]  $ \log(\sqrt[ 4 ]{{0.03}^3}) $  \begin{solution}  $ -1.14215906 $  \end{solution} \part[1]  $ \log(\sqrt[ 5 ]{0.81}) $  \begin{solution}  $ -0.0183029962 $  \end{solution} \part[1]  $ \log(\frac{1}{81}) $  \begin{solution}  $ -1.90848502 $  \end{solution}
        \end{parts}
        \end{multicols}
        \question p028e11 - Sabiendo que log 2 = 0,301030 y que log 3 = 0,477121 averigua, sin calculadora:
        \begin{multicols}{3} 
        \begin{parts} \part[1]  $ \log(8) $  \begin{solution}  $ 0.903089987 $  \end{solution} \part[1]  $ \log(15) $  \begin{solution}  $ 1.17609126 $  \end{solution} \part[1]  $ \log(12) $  \begin{solution}  $ 1.07918125 $  \end{solution} \part[1]  $ \log(\sqrt[ 3 ]{162}) $  \begin{solution}  $ 0.736505005 $  \end{solution} \part[1]  $ \log(0.18) $  \begin{solution}  $ -0.744727495 $  \end{solution} \part[1]  $ \log(0.002) $  \begin{solution}  $ -2.69897 $  \end{solution} \part[1]  $ \log(0.0625) $  \begin{solution}  $ -1.20411998 $  \end{solution} \part[1]  $ \log(40.5) $  \begin{solution}  $ 1.60745502 $  \end{solution} \part[1]  $ \log(\frac{1}{3}) $  \begin{solution}  $ -0.477121255 $  \end{solution}
        \end{parts}
        \end{multicols}
        \question p028e12 - Sabiendo que log 2 = 0,301030 y que log 3 = 0,477121 averigua, sin calculadora:
        \begin{multicols}{3} 
        \begin{parts} \part[1]  $ \log \frac{{{{0.0027}^3} \cdot \sqrt[ 4 ]{540} }}{{96 \cdot \sqrt[ 5]{51.84} }} $  \begin{solution}  $ -9.3480145 $  \end{solution}
        \end{parts}
        \end{multicols}
        \question p029e15 - Realiza las siguientes operaciones:
        \begin{multicols}{3} 
        \begin{parts} \part[1]  $ \log ( {4 - \sqrt{ 6} } ) + \log ( {4 + \sqrt {6} } ) $  \begin{solution}  $ 1 $  \end{solution} \part[1]  $ \frac{1}{2}\log ( {12 - 2\sqrt {11} } ) + \frac{1}{2}\log ( {12 + 2\sqrt {11} } ) $  \begin{solution}  $ 1 $  \end{solution} \part[1]  $ \log ( {7 - \sqrt {22} } ) + \log ( {7 + \sqrt {22} } ) - 3\log 3 $  \begin{solution}  $ 0 $  \end{solution}
        \end{parts}
        \end{multicols}
        \question p029e18 - Averigua el valor de x en los siguientes casos:
        \begin{multicols}{2} 
        \begin{parts} \part[1]  $ \log (x) - \log (3) = 2 $  \begin{solution}  $ \left [ 300\right ] $  \end{solution} \part[1]  $  \log 2 + \log x = 1 $  \begin{solution}  $ \left [ 5\right ] $  \end{solution} \part[1]  $ 5\log x - \log 32 = \log(\frac{x}{2}) $  \begin{solution}  $ \left [ 2\right ] $  \end{solution} \part[1]  $ 2\log x - \log (x - 16) = 2 $  \begin{solution}  $ \left [ \right ] $  \end{solution} \part[1]  $ \log x + \log (50) = \log (1000) $  \begin{solution}  $ \left [ 20\right ] $  \end{solution} \part[1]  $ \log (x^3) = \log (6) + 2 \log (x) $  \begin{solution}  $ \left [ 6\right ] $  \end{solution} \part[1]  $ 2\log (\frac{x}{2}) + 2\log (\frac{x}{3}) = 3\log (x) - \log (\frac{{32}}{9}) $  \begin{solution}  $ \left [ \frac{81}{8}\right ] $  \end{solution} \part[1]  $ \log(2) + \log(11-x^2) = 2\log(5-x) $  \begin{solution}  $ \left [ \frac{1}{3}, \quad 3\right ] $  \end{solution} \part[1]  $ \log(1250) - 2 = 2 - \log {( {{2^{2 - x}}} )^{2 + x}} $  \begin{solution}  $ \left [ -1, \quad 1\right ] $  \end{solution} \part[1]  $ \log(x - 1) - \log (\sqrt{5 + x})  - \log(\sqrt{5 - x})=0 $  \begin{solution}  $ \left [ 4\right ] $  \end{solution} \part[1]  $ {\log (x - 1) + \log (x + 1)} = 2{\log (2 - x)} $  \begin{solution}  $ \left [ \frac{5}{4}\right ] $  \end{solution}
        \end{parts}
        \end{multicols}
        \question p029e19 - Resuelve los siguientes sistemas de inecuaciones:
        \begin{multicols}{2} 
        \begin{parts} \part[1]  $ \left\{\begin{matrix}\log x + \log y = 5 \\\log x - \log y = 1\end{matrix}\right. $  \begin{solution}  $ \left [ \left \{ x : 1000, \quad y : 100\right \}\right ] $  \end{solution} \part[1]  $ \left\{\begin{matrix}3\log x - 2\log y = 10\\\log x + 3\log y = 7\end{matrix}\right. $  \begin{solution}  $ \left [ \left \{ x : 10000, \quad y : 10\right \}\right ] $  \end{solution} \part[1]  $ \left\{\begin{matrix}\log_{2}(x) + \log_{2}(x + y) = 4\\x + y = 2 \end{matrix}\right. $  \begin{solution}  $ \left [ \left \{ x : 8, \quad y : -6\right \}\right ] $  \end{solution}
        \end{parts}
        \end{multicols}
        
    \end{questions}
    \end{document}
    